% Companion manuscript. Compile from the repository root.
% Keep main_rewrite.tex and its finite-frontier results independent.
\documentclass[11pt]{article}
\usepackage[margin=1in]{geometry}
\usepackage{amsmath,amssymb,amsthm,mathtools}
\usepackage{booktabs,array,float}
\usepackage{microtype,setspace,needspace}
\usepackage[round,authoryear]{natbib}
\usepackage[colorlinks=true,allcolors=blue]{hyperref}
\allowdisplaybreaks[1]
\newtheorem{proposition}{Proposition}
\newtheorem{definition}{Definition}
\newcommand{\dd}{\,\mathrm{d}}
\newcolumntype{P}[1]{>{\raggedright\arraybackslash}p{#1}}

\title{How Much Can Recursive AI Self-Improvement\\Raise Economic Growth?}
\author{Oliver Pardo}
\date{September 2026}

\begin{document}
\maketitle
\begin{abstract}
I study how recursive AI self-improvement affects long-run economic growth
in a decentralized economy with autonomous AI research. A monopolistic
developer supplies AI services and invests in AI efficiency without a finite
efficiency frontier. With unit elasticity of substitution between AI services
and labor, I derive a balanced-growth equilibrium in which research spending,
saving, and prices are endogenous. The increase in long-run growth depends on
research returns, the production weight of AI, and exogenous effective-labor
growth. Output per person and real wages grow at the same rate, while labor
retains a positive income share. I use this analytical benchmark to organize
calibration and parameter sensitivity. The inherited illustrative calibration
implies a growth increase of about 0.087 percentage points per year; this is
a conditional model calculation, not an empirical estimate or a forecast.
\end{abstract}
\onehalfspacing
\section{Introduction}
\label{sec:companion-introduction}

How much can economic growth increase if AI systems help develop better AI?
Answering this question requires connecting the productivity of AI research
to the allocation of resources in the rest of the economy. Research competes
with current consumption, physical investment, and the compute used to supply
AI services. Its scale depends on the profits that improvements can earn and
on the interest rate at which those profits are discounted.

I study these decisions in an endogenous-growth model with a representative
household, competitive final-good firms, and a monopolistic AI developer.
Final production combines capital, labor, and AI services. The developer
uses AI research services to raise the efficiency with which compute produces
both production and research services. Research is autonomous: the model
describes an economy in which human researchers are no longer a necessary
current input into AI development. It does not describe how that transition
to autonomous research occurs.

This paper develops the uncapped unit-elastic benchmark in Appendix C of
\citet{pardo2026future}. That paper compares substitution regimes under a
finite AI-efficiency frontier. Here I fix the elasticity of substitution at
one and remove that frontier. I derive an equilibrium with constant growth
rates and endogenous allocation shares, rather than assume fixed saving and
research expenditure. This restriction isolates a tractable quantitative
question; it is not a claim that aggregate AI--labor substitution has been
estimated to equal one.

I distinguish inherited macroeconomic reference values from uncertain AI
parameters. The analytical results identify which parameters affect permanent
growth and which affect reference levels. This first draft establishes that
benchmark and the calibration framework. Empirically disciplined ranges for
the AI parameters and quantitative sensitivity exercises remain to be developed.
The calculations do not assume that the present world economy lies on the
model's balanced-growth path.

\section{Model}
\label{sec:companion-model}

Time is continuous and measured in years. Population $N$ and labor-augmenting
technology $A$ satisfy
\begin{equation}
 \dot N=nN,\qquad \dot A=\gamma A,
 \qquad n\geq0,\quad\gamma>0.
 \label{eq:companion-exogenous}
\end{equation}
Each person supplies one unit of labor inelastically. Production labor is
denoted by $L$; labor-market clearing will impose $L=N$. For any positive
variable $V$, write $g_V=\dot V/V$. The final good is the numeraire.

\subsection{Household}

The household owns physical capital $K$ and the AI developer. Let $C$ denote
aggregate consumption, $w$ the real wage, $r$ the net real interest rate, and
$\Pi$ the developer's net distribution to its owner. Given their paths and
$K_0>0$, the household chooses consumption to solve
\begin{equation}
 \max\int_0^\infty e^{-\rho t}N_t\log(C_t/N_t)\dd t,
 \qquad \rho>n,
 \label{eq:companion-household-objective}
\end{equation}
subject to
\begin{equation}
 \dot K=rK+wN+\Pi-C,\qquad C>0,\quad K\geq0.
 \label{eq:companion-household-budget}
\end{equation}
The parameter $\rho$ is the discount rate. For an interior optimum, the Euler
and transversality conditions are
\begin{align}
 g_C&=n+r-\rho,\label{eq:companion-euler}\\
 \lim_{t\to\infty}e^{-\rho t}\frac{N_tK_t}{C_t}&=0.
 \label{eq:companion-household-tvc}
\end{align}
Depreciation is already deducted from $r$.

\subsection{Final-good firms}

Let $X$ denote AI production services. I set the elasticity of substitution
between effective labor $AL$ and $X$ to $\sigma=1$, retaining the production
weights $\omega_X\in(0,1)$ and $\omega_L=1-\omega_X$. With capital elasticity
$\alpha\in(0,1)$, production is
\begin{equation}
 Y=K^\alpha\left[(AL)^{\omega_L}X^{\omega_X}\right]^{1-\alpha}.
 \label{eq:companion-production-original}
\end{equation}
Define the output elasticities
\begin{equation}
 \beta=(1-\alpha)\omega_X,\qquad
 \lambda=(1-\alpha)\omega_L,\qquad\alpha+\beta+\lambda=1.
 \label{eq:companion-elasticities}
\end{equation}
Thus the same production function is $Y=K^\alpha(AL)^\lambda X^\beta$.
Let $p_X$ denote the AI-service price and $\delta\geq0$ capital depreciation.
A competitive final-good firm solves
\begin{equation}
 \max_{K,L,X\geq0}\{K^\alpha(AL)^\lambda X^\beta-(r+\delta)K-wL-p_XX\}.
 \label{eq:companion-firm}
\end{equation}
Its interior first-order conditions are
\begin{align}
 r&=\alpha Y/K-\delta,\label{eq:companion-rental}\\
 w&=\lambda Y/L,\label{eq:companion-wage}\\
 p_X&=\beta Y/X.\label{eq:companion-demand}
\end{align}
Constant returns imply zero final-good profit. At given $K,A,L$, equation
\eqref{eq:companion-demand} defines inverse demand $p_X(X)$ for the developer.

\subsection{The AI developer}

Inference compute $U$ and research compute $M$ are flows of computational
resources, each costing one unit of the final good per unit of compute.
They are rival inputs, not stocks of processors. AI efficiency $B$ converts
them into the flows $BU$ of AI production services and $BM$ of AI research
services. The same nonrival efficiency index augments both uses of compute.
The service technology and research law are
\begin{align}
 X&=BU,\label{eq:companion-services}\\
 \dot B&=\chi(BM)^\eta,\qquad\chi>0,\quad0<\eta<1.
 \label{eq:companion-research-law}
\end{align}
The parameter $\chi$ shifts research productivity; $\eta$ is the elasticity
of the efficiency increment $\dot B$ with respect to research services $BM$.
Efficiency has no finite upper bound. Higher $B$ raises research productivity
at a given $M$, creating recursive self-improvement. There is no human
research input in this benchmark.

The developer exclusively owns $B$, sells $X$, and purchases $U=X/B$ and $M$.
Ownership makes some benefits of research appropriable through monopoly rents.
Taking $K,A,L,r$ as given and starting from $B_0>0$, it maximizes
\begin{equation}
 \int_0^\infty D_t[p_X(X_t)X_t-X_t/B_t-M_t]\dd t,
 \qquad D_t=\exp\!\left(-\int_0^t r_s\dd s\right),
 \label{eq:companion-developer-objective}
\end{equation}
subject to \eqref{eq:companion-research-law}, $X,M\geq0$, and finite feasible
allocations at every finite date. Since $X$ does not enter the research law
directly, define optimized operating profit, before research expenditure, by
\begin{equation}
 \pi(B)=\max_{X\geq0}\{p_X(X)X-X/B\}.
 \label{eq:companion-static-monopoly}
\end{equation}
Dependence on current $K,A,L$ is implicit. Revenue is proportional to $X^\beta$,
so the objective is strictly concave, its derivative is positive near zero,
and linear cost dominates revenue at infinity. The unique positive optimum
therefore satisfies
\begin{equation}
 \beta p_X=1/B.
 \label{eq:companion-monopoly-foc}
\end{equation}
Combining this condition with \eqref{eq:companion-demand} and
\eqref{eq:companion-services}, and applying the envelope theorem, gives
\begin{equation}
 p_XX=\beta Y,\qquad U=\beta^2Y,\qquad\pi_B=U/B.
 \label{eq:companion-static-shares}
\end{equation}

Let $q$ be the current-value shadow price of AI efficiency. The Hamiltonian
is $\pi(B)-M+q\chi(BM)^\eta$. Interior research and the costate condition give
\begin{align}
 1&=q\chi\eta B^\eta M^{\eta-1},\label{eq:companion-research-foc}\\
 \dot q&=rq-U/B-\eta q\dot B/B.\label{eq:companion-costate}
\end{align}
Research equates its unit resource cost to the value of the efficiency it
creates. Holding efficiency also yields compute-cost savings $U/B$ and
raises the productivity of subsequent research. The required return $rq$
equals those two benefits plus the capital gain $\dot q$. The developer's
transversality condition is
\begin{equation}
 \lim_{t\to\infty}D_tq_tB_t=0.
 \label{eq:companion-developer-tvc}
\end{equation}
Net distributions are $\Pi=\pi(B)-M=p_XX-U-M$ and can be negative.

\subsection{Equilibrium}

\begin{definition}
\label{def:companion-equilibrium}
Given parameters and initial conditions $A_0,N_0,K_0,B_0>0$, an interior
equilibrium is a trajectory for $(K,B,C,q,Y,L,X,U,M,r,w,p_X,\Pi)$, finite at
each finite date, with $K,B,C,q,Y,L,X,U,M>0$, such that the household and
developer attain their global optima, final-good firms maximize profit, and
markets clear. Its equations and boundary conditions are
\begin{align}
 A_t&=A_0e^{\gamma t},\qquad N_t=N_0e^{nt},\qquad L_t=N_t,
 \label{eq:companion-eq-exogenous}\\
 Y&=K^\alpha(AL)^\lambda X^\beta,\label{eq:companion-eq-output}\\
 X&=BU,\label{eq:companion-eq-services}\\
 p_X&=\beta Y/X,\qquad\beta p_X=1/B,\label{eq:companion-eq-price}\\
 r&=\alpha Y/K-\delta,\qquad w=\lambda Y/L,\label{eq:companion-eq-factors}\\
 \Pi&=p_XX-U-M,\label{eq:companion-eq-profit}\\
 \dot K&=Y-C-U-M-\delta K,\label{eq:companion-eq-kdot}\\
 \dot B&=\chi(BM)^\eta,\label{eq:companion-eq-bdot}\\
 \dot C&=(n+r-\rho)C,\label{eq:companion-eq-cdot}\\
 1&=q\chi\eta B^\eta M^{\eta-1},\label{eq:companion-eq-research}\\
 \dot q&=rq-U/B-\eta q\dot B/B,\label{eq:companion-eq-qdot}\\
 K(0)&=K_0,\qquad B(0)=B_0,\label{eq:companion-eq-initial}\\
 \lim_{t\to\infty}e^{-\rho t}N_tK_t/C_t&=0,\label{eq:companion-eq-household-tvc}\\
 \lim_{t\to\infty}D_tq_tB_t&=0.\label{eq:companion-eq-developer-tvc}
\end{align}
\end{definition}
Capital accumulation is the aggregate resource constraint; the household
budget follows from it and the factor-payment and distribution equations.
The initial jump variables $C_0,q_0$ are endogenous. This definition imposes
no balanced-growth restrictions. The first-order conditions and TVCs alone
are not asserted to be sufficient for equilibrium: global optimality is
verified under explicit conditions below.

\section{Balanced growth and equilibrium existence}
\label{sec:companion-bgp}

Substituting $X=B\beta^2Y$ from \eqref{eq:companion-static-shares} into
\eqref{eq:companion-production-original} yields
\begin{equation}
 Y^{1-\beta}=K^\alpha(AL)^\lambda(\beta^2B)^\beta.
 \label{eq:companion-reduced-output}
\end{equation}
Consider a candidate balanced-growth path (BGP) with positive constant
$K/Y,C/Y,M/Y$. Differentiating \eqref{eq:companion-reduced-output} and the
research law \eqref{eq:companion-research-law} gives, respectively,
\begin{align}
 g_Y&=n+\gamma+(\beta/\lambda)g_B,\label{eq:companion-production-growth}\\
 (1-\eta)g_B&=\eta g_Y.\label{eq:companion-research-growth}
\end{align}
The first equation relates AI improvement to aggregate production. The second
requires research compute to grow with output when its expenditure share is
constant. Define
\begin{equation}
 \Delta=\lambda(1-\eta)-\beta\eta=(1-\alpha)(1-\eta-\omega_X).
 \label{eq:companion-feedback}
\end{equation}
If $\Delta>0$, the two growth equations have the solution
\begin{align}
 g_Y^*&=\frac{(1-\eta)(1-\omega_X)}{1-\eta-\omega_X}(n+\gamma),
 \label{eq:companion-gy}\\
 g_B^*&=\frac{\eta(1-\omega_X)}{1-\eta-\omega_X}(n+\gamma).
 \label{eq:companion-gb}
\end{align}
The superscript $*$ identifies this BGP. The wage equation and household Euler
condition then imply
\begin{align}
 g_{Y/N}^*=g_w^*&=\gamma+
 \frac{\eta\omega_X}{1-\eta-\omega_X}(n+\gamma),
 \label{eq:companion-premium}\\
 r^*&=\rho+g_{Y/N}^*.\label{eq:companion-bgp-return}
\end{align}

To check that these rates belong to an equilibrium, define the candidate
ratios $k^*=K/Y$, $u^*=U/Y$, $i^*=(\dot K+\delta K)/Y$, $m^*=M/Y$, and
$c^*=C/Y$. These are capital--output and expenditure ratios, not quantities
per unit of effective labor. Their values are
\begin{align}
 k^*&=\frac{\alpha}{r^*+\delta},\qquad u^*=\beta^2,
 \label{eq:companion-bgp-capital-inference}\\
 i^*&=(g_Y^*+\delta)k^*,\label{eq:companion-bgp-investment}\\
 m^*&=\frac{\beta^2\eta g_B^*}{\rho-n+\eta g_Y^*},
 \label{eq:companion-bgp-research}\\
 c^*&=1-u^*-i^*-m^*.\label{eq:companion-bgp-consumption}
\end{align}
The research ratio follows from the developer's conditions, not a fixed
research rule. Indeed, \eqref{eq:companion-research-foc} implies
$M=\eta qg_BB$. Constant $M/Y$ then gives $g_q=g_Y-g_B$. Dividing
\eqref{eq:companion-costate} by $q$ and substituting these expressions yields
\eqref{eq:companion-bgp-research}. The remaining ratios follow from the
capital return, the static monopoly solution, and resource clearing.

\Needspace{7\baselineskip}
\begin{proposition}[A balanced-growth equilibrium]
\label{prop:companion-bgp}
Maintain the primitive restrictions in Section~\ref{sec:companion-model}.
Suppose, in addition, that
\begin{equation}
 0<\eta<\alpha,\qquad\eta\leq\tfrac12,\qquad
 \beta\leq\tfrac12,\qquad\eta+\omega_X<1.
 \label{eq:companion-sufficient}
\end{equation}
For every $A_0,N_0>0$, there exist $K_0^*,B_0^*>0$ and an equilibrium with
the rates and ratios in \eqref{eq:companion-gy}--\eqref{eq:companion-bgp-consumption}.
Consumption is positive, both agents attain finite-valued global optima,
and both transversality conditions hold. AI efficiency grows without bound
as time tends to infinity but remains finite at every finite date.
\end{proposition}
\begin{proof}
See Appendix~\ref{app:companion-proofs}, Proof of
Proposition~\ref{prop:companion-bgp}.
\end{proof}

The restrictions serve different purposes. The condition $\eta+\omega_X<1$
is necessary for this positive finite-rate BGP when $n+\gamma>0$.
If it fails, the two growth restrictions either contradict one another or
require declining AI efficiency. This excludes the specified BGP, not every
possible equilibrium trajectory. The bounds $\beta\leq1/2$ and $\eta\leq1/2$
are sufficient concavity conditions for global developer optimality. They
are not claimed to be necessary. I retain $\eta<\alpha$ for comparability
with the companion framework in \citet{pardo2026future}; the proof below
does not separately require that comparison.

Proposition~\ref{prop:companion-bgp} constructs the initial stocks together
with the BGP. It does not prove equilibrium existence for arbitrary fixed
$K_0,B_0$. Appendix C of \citet{pardo2026future} also establishes local
equilibrium trajectories from nearby initial stocks under explicit stability
and state-projection conditions. No off-BGP simulation is presented here.

The BGP income shares are
\begin{equation}
 \frac{wL}{Y}=\lambda,\qquad
 \frac{(r+\delta)K}{Y}=\alpha,\qquad
 \frac{p_XX}{Y}=\beta,\qquad
 \frac{\Pi}{Y}=\beta-\beta^2-m^*.
 \label{eq:companion-distribution}
\end{equation}
Real wages and output per person grow together faster than $\gamma$ while
labor retains a positive income share. This differs from the characterized
unit-elastic finite-frontier equilibria, where both growth rates tend to
$\gamma$. Removing the frontier changes the long-run result; the uncapped
BGP is not obtained by interchanging the long-run and infinite-frontier limits.

\section{Calibration and measurement}
\label{sec:companion-calibration}

Table~\ref{tab:companion-parameters} imports the reference values documented
in the quantitative section of \citet{pardo2026future}. It separates
macroeconomic benchmarks from illustrative AI parameters. Reusing these
values preserves comparability; it does not make them empirical estimates
of this autonomous-research economy.

\begin{table}[H]
\centering
\caption{Inherited parameters and their evidential status}
\label{tab:companion-parameters}
\small
\setlength{\tabcolsep}{4pt}
\renewcommand{\arraystretch}{1.12}
\begin{tabular}{P{0.10\textwidth}P{0.11\textwidth}P{0.26\textwidth}P{0.42\textwidth}}
\toprule
Parameter & Reference & Interpretation & Source or status\\
\midrule
$\alpha$ & 0.33 & Capital elasticity & Inherited one-third macroeconomic benchmark.\\
$\delta$ & 0.05 & Annual depreciation & Inherited benchmark; PWT is the comparison source, not a direct estimate of this value \citep{pwt110}.\\
$\rho$ & 0.04 & Annual discount rate & Preference calibration, not a measured market interest rate.\\
$n$ & 0.003 & Annual population growth & Inherited constant-rate approximation to the UN 2024--2100 world projection \citep{unwpp2024}.\\
$\gamma$ & 0.01 & Annual labor-augmenting progress & Illustrative trend, with the OECD long-run scenarios as a scale comparison \citep{oecdlongrun2025}.\\
$\omega_X$ & 0.20 & AI production weight & Illustrative, not estimated. Not the AI revenue share itself.\\
$\omega_L$ & 0.80 & Effective-labor weight & Defined by $1-\omega_X$.\\
$\eta$ & 0.20 & Research-services elasticity & Illustrative, not estimated.\\
$\sigma$ & 1 & AI--labor substitution & Maintained restriction of this paper.\\
$\chi$ & Free, $>0$ & Research productivity & Does not affect the BGP rates or expenditure ratios. Needs units and level targets for a quantitative level comparison.\\
\bottomrule
\end{tabular}
\end{table}

The population parameter extrapolates a finite-period approximation, not a
UN prediction of positive population growth forever. Likewise, measured
labor-productivity growth is not the same object as the growth of $A$ in
this model. The OECD comparison does not directly identify $\gamma$.
Both $n$ and $\gamma$ therefore belong in long-run sensitivity analysis.

The finite-frontier paper uses $\chi=1.4378$ to target an illustrative
transition duration. I do not transfer that timing target to the uncapped
economy. Equations~\eqref{eq:companion-gy}--\eqref{eq:companion-bgp-consumption}
are independent of $\chi$, whereas the initial BGP levels constructed in
Appendix~\ref{app:companion-proofs} depend on it. Calibrating $\chi$ cannot
select an arbitrary permanent growth rate.

At the inherited reference values, the analytical calculations give
\begin{equation}
 g_{Y/N}^*=g_w^*\simeq1.0867\%\ \text{per year},\qquad
 r^*\simeq5.0867\%\ \text{per year}.
 \label{eq:companion-reference-results}
\end{equation}
The growth increase relative to $\gamma=1\%$ is about $0.0867$ percentage
points per year. This comparison is between long-run model benchmarks, not
a transition from today's economy and not an estimate of AI's observed impact.

Empirical discipline requires a mapping from observables to model objects.
At unit elasticity, \eqref{eq:companion-distribution} implies that a measured
AI revenue share would map to $\omega_X$ only after dividing by $1-\alpha$,
provided the measurement matches the model's industry boundary and pricing
assumptions. AI investment, compute expenditure, and AI revenue are different
objects. Measurement difficulties discussed by \citet{oecdaiinvestment2025}
motivate caution; they do not select $\omega_X=0.20$.

For research, taking logs of \eqref{eq:companion-research-law} gives
\begin{equation}
 \log\dot B=\log\chi+\eta\log B+\eta\log M.
 \label{eq:companion-measurement-law}
\end{equation}
The left side is the log of an efficiency increment, not the growth rate
$g_B$. Moreover, research compute is an endogenous choice. Historical
algorithmic improvements cannot be read directly as an estimate of $\eta$:
they may also reflect human research, data, and complementary inputs omitted
from the autonomous benchmark. Identifying defensible ranges for $\eta$ and
$\omega_X$ remains a separate empirical task.

\section{Sensitivity and the quantitative question}
\label{sec:companion-sensitivity}

The analytical sensitivity follows directly from
\eqref{eq:companion-premium}. Within the maintained existence region,
\begin{align}
 \frac{\partial g_{Y/N}^*}{\partial\eta}
 &=\frac{\omega_X(1-\omega_X)}{(1-\eta-\omega_X)^2}(n+\gamma)>0,
 \label{eq:companion-sensitivity-eta}\\
 \frac{\partial g_{Y/N}^*}{\partial\omega_X}
 &=\frac{\eta(1-\eta)}{(1-\eta-\omega_X)^2}(n+\gamma)>0.
 \label{eq:companion-sensitivity-omega}
\end{align}
The same derivatives apply to real-wage growth and to the interest rate
when $\rho$ is fixed. These are comparisons of equilibria, not partial
adjustments holding research expenditure fixed.

The growth increase is amplified as $\eta+\omega_X$ approaches one from
below. This makes the empirical plausibility of the parameter combinations
particularly important. Such sensitivity is not itself evidence that large
growth increases are likely. Parameter values outside the sufficient
conditions of Proposition~\ref{prop:companion-bgp} must not be presented as
equilibria merely because the closed-form growth expression can be evaluated.

The next quantitative exercise will vary the two AI parameters jointly and
report output-per-person growth, wage growth, the interest rate, and the
research expenditure share. Every configuration must satisfy the stated
existence restrictions. Sensitivity to $n$ and $\gamma$ will distinguish
assumptions about AI from assumptions about exogenous effective-labor growth.
Changes in $\alpha,\delta,\rho$ also matter for allocations: their absence
from the growth premium does not make them irrelevant to the equilibrium.

The no-AI growth benchmark is $\gamma$. Comparing the positive-AI BGP with
that benchmark combines the production and recursive-research channels.
It does not isolate the causal contribution of recursion. A counterfactual
that retains AI services in production while removing their contribution to
research productivity requires a separately specified and solved research
technology; it is not implemented in this draft.

The current contribution is an analytical, globally verified BGP benchmark
and a transparent starting calibration. The magnitude implied by the
reference values is not a forecast. The next step is to discipline the
uncertain AI parameters with evidence and then report conditional ranges,
without treating sensitivity scenarios as estimated probabilities.

\appendix
\section{Proofs}
\label{app:companion-proofs}

\begin{proof}[Proof of Proposition~\ref{prop:companion-bgp}]
First, all candidate ratios are positive. The denominator in
\eqref{eq:companion-bgp-research} is positive, and
\[
 i^*=\frac{\alpha(g_Y^*+\delta)}{\rho-n+g_Y^*+\delta}<\alpha,
 \qquad m^*<\frac{\beta^2\eta}{1-\eta}.
\]
Since $\Delta=(1-\alpha)(1-\eta)-\beta>0$, both
$\beta<1-\eta$ and $\beta<1-\alpha$. Hence
\[
 u^*+m^*<\frac{\beta^2}{1-\eta}<\beta<1-\alpha,
 \qquad c^*=1-u^*-i^*-m^*>0.
\]
Positive consumption is a consequence, not an additional assumption.

For given $A_0,N_0>0$, construct initial levels as
\begin{align}
 B_0^*&=\left[A_0N_0(k^*)^{\alpha/\lambda}(u^*)^{\beta/\lambda}
 m^*\left(\frac{\chi}{g_B^*}\right)^{1/\eta}
 \right]^{\eta\lambda/\Delta},\label{eq:companion-initial-b}\\
 Y_0^*&=A_0N_0(k^*)^{\alpha/\lambda}(u^*B_0^*)^{\beta/\lambda},
 \qquad K_0^*=k^*Y_0^*,\label{eq:companion-initial-y}\\
 q_0^*&=\frac{m^*Y_0^*}{\eta g_B^*B_0^*}.
 \label{eq:companion-initial-q}
\end{align}
The first two equations solve the initial production equation and
$g_B^*B_0^*=\chi(B_0^*m^*Y_0^*)^\eta$ simultaneously.
Let $Y_t^*=Y_0^*e^{g_Y^*t}$, $B_t^*=B_0^*e^{g_B^*t}$, and
$q_t^*=q_0^*e^{(g_Y^*-g_B^*)t}$. Set
$K=k^*Y^*$, $C=c^*Y^*$, $U=u^*Y^*$, $M=m^*Y^*$, $X=B^*U$, and $L=N$.
Determine $p_X,r,w,\Pi$ from the static equations in
Definition~\ref{def:companion-equilibrium}.

Equations~\eqref{eq:companion-production-growth} and
\eqref{eq:companion-research-growth} make production and the research law
hold at every date. The chosen $q_0^*$ ensures $M=\eta q\dot B$, which is
the research FOC. The capital-return, Euler, and resource equations follow
from $k^*,r^*,c^*$. Finally, the costate equation reduces to
\[
 g_Y^*-g_B^*=r^*-\frac{u^*\eta g_B^*}{m^*}-\eta g_B^*,
\]
which follows from \eqref{eq:companion-bgp-research} and
$(1-\eta)g_B^*=\eta g_Y^*$. All quantities remain finite and positive at
every finite date.

Both transversality conditions hold explicitly:
\begin{align}
 e^{-\rho t}\frac{N_tK_t^*}{C_t^*}
 &=\frac{N_0k^*}{c^*}e^{-(\rho-n)t}\longrightarrow0,
 \label{eq:companion-proof-household-tvc}\\
 D_tq_t^*B_t^*&=q_0^*B_0^*e^{-(\rho-n)t}\longrightarrow0.
 \label{eq:companion-proof-developer-tvc}
\end{align}
Here $D_t=e^{-r^*t}$. Discounted developer profit is proportional to
$e^{-(\rho-n)t}$. Discounted household utility is
$N_0e^{-(\rho-n)t}[\log(C_0^*/N_0)+(g_Y^*-n)t]$. Both objective integrals
are finite because $\rho>n$.

To verify global developer optimality, static optimization at given
$K_t,A_t,L_t$ gives
\begin{equation}
 \pi_t(B)=\beta(1-\beta)
 [K_t^\alpha(A_tL_t)^\lambda]^{1/(1-\beta)}
 (\beta^2B)^{\beta/(1-\beta)}.
 \label{eq:companion-profit-function}
\end{equation}
This function is concave in $B$ when $\beta\leq1/2$. Research technology
$F(B,M)=\chi B^\eta M^\eta$ is jointly concave when $\eta\leq1/2$.
Since $q_t>0$, the reduced Hamiltonian $\pi_t(B)-M+q_tF(B,M)$ is jointly
concave. Let $J_T=\int_0^T D_t[\pi_t(B_t)-M_t]\dd t$.
For any admissible alternative, denoted by a tilde, with the same $B_0$,
the supporting-hyperplane inequality, research FOC, and costate equation give
\begin{equation}
 \widetilde J_T-J_T
 \leq-D_Tq_T(\widetilde B_T-B_T)
 \leq D_Tq_TB_T\longrightarrow0.
 \label{eq:companion-developer-verification}
\end{equation}
No alternative can have a larger limiting payoff, or a larger upper limit
of truncated payoff, than the constructed finite-valued candidate. Static
optimality was global by strict concavity of
\eqref{eq:companion-static-monopoly}, so reducing the developer problem did
not discard a profitable service choice.

For the household, write $\Lambda_t=e^{-\rho t}N_t/C_t$. Euler implies
$\dot\Lambda=-r\Lambda$. Concavity of log utility and the linear household
budget yield, for any feasible alternative at the same given prices and
distributions and with the same $K_0$,
\begin{equation}
 \int_0^T e^{-\rho t}N_t
 [\log(\widetilde C_t/N_t)-\log(C_t/N_t)]\dd t
 \leq-\Lambda_T(\widetilde K_T-K_T)
 \leq\Lambda_TK_T\longrightarrow0.
 \label{eq:companion-household-verification}
\end{equation}
This proves global household optimality. Final-good firms face a concave
constant-returns technology, so their first-order conditions maximize profit.
Market clearing and both global agent optima establish the claimed equilibrium.
\end{proof}

\bibliographystyle{apalike}
\bibliography{references}
\end{document}
